\section{工程化视角：Q-Learning 工作流}

\subsection{数据→训练→部署闭环}

工程项目的典型闭环包含：数据收集 → 模型训练 → 评估复核 → Sim rollout + human feedback → 线上监控。

\begin{importantbox}{闭环流程关键节点}
1. 收集 high-quality transitions，避免 agent 作弊。2. 用 PER + Double Q + regularization 做训练稳定化。3. 增强评估范式：human-in-the-loop replay、failure case tagging。4. 设定 monitoring guardrail：TD error、reward variance、latency。
\end{importantbox}

\subsection{部署监控策略}

部署后重点监控：

\begin{itemize}
    \item `TD error` drift（若上升需重新训练/rollback）
    \item reward variance（如变小可能代表 behavior collapse）
    \item `episode success rate` 与 `reward per step`，确保性能不退化
\end{itemize}

\subsection{本章小结}

工程化部署强调 replay buffer hygiene、TD error 警戒和 telemetry-driven rollback。

